\chapter{Tight-Binding and Slater--Koster Representations}
\index{tight-binding model|textbf}
\index{Slater--Koster representation|textbf}

\label{ch:tight-binding-slater-koster}

\chapterstatus{mathematical and scientific definition of the orthogonal,
spinless, homopolar $sp^3s^{\ast}$ representation used by the bulk-silicon
reduction program. This chapter defines a model class; it does not report a
fit, accepted parameter set, or validated silicon Hamiltonian.}

A tight-binding Hamiltonian is a represented operator on a declared localized
coordinate space. A Slater--Koster model further restricts its two-center
matrix elements by orbital character, bond direction, parameter sharing, and
interaction range~\cite{slater1954,vogl1983}. These restrictions are physical
model assumptions, not merely an efficient storage scheme. They must therefore
be fixed before a fit is interpreted.

This chapter separates four objects that are often conflated:

\begin{enumerate}
  \item the ordered localized basis and its orbital centers;
  \item the ordinary-cell hopping coefficients and Bloch convention;
  \item the two-center Slater--Koster angular construction; and
  \item the material-specific bond inventory, onsite terms, and parameter
        sharing that compose a candidate Hamiltonian.
\end{enumerate}

Parameter estimation, residual analysis, admissible sets, withheld validation,
and scientific acceptance remain the subjects of
Chapter~\ref{ch:model-reduction}.

\section{Localized coordinates and represented state space}
\index{basis!tight-binding}
\index{orbital center}

Let $\mathcal L\cong\mathbb Z^3$ be the ordinary-cell translation lattice and
let
\begin{equation}
  \mathcal C=(c_0,c_1,\ldots,c_{N_c-1})
\end{equation}
be an ordered inventory of centers in one primitive cell. Center $c$ has a
fractional coordinate
\begin{equation}
  \boldsymbol\tau_c\in\mathbb R^3.
\end{equation}
If $A\in\mathbb R^{3\times3}$ contains the direct-lattice vectors as columns,
then $A(\mathbf R+\boldsymbol\tau_c)$ is the Cartesian position of center $c$
in ordinary cell $\mathbf R\in\mathbb Z^3$.

A localized spatial basis coordinate is the triple
\begin{equation}
  b_a=(c_a,\lambda_a,r_a),
\end{equation}
where $c_a$ is a center, $\lambda_a$ is a real-cubic orbital component, and
$r_a\in\mathbb N_0$ is a radial-channel label. The represented spatial state
space is
\begin{equation}
  \mathcal H_{\mathrm{TB}}
  \cong
  \ell^2(\mathcal L)\otimes\mathbb C^{N_b}.
\end{equation}
The coordinate labels define a finite representation. They do not by
themselves construct continuum atomic orbitals, prove completeness, or validate
a material parameterization.

\subsection{The homopolar \texorpdfstring{$sp^3s^{\ast}$}{sp3s*} frame}
\index{tight-binding model!sp3s-star}

For the admitted homopolar $sp^3s^{\ast}$ class, each center contributes the
ordered five-coordinate block
\begin{equation}
  \mathcal B_c
  =
  \bigl(
    (c,s,0),
    (c,p_x,0),
    (c,p_y,0),
    (c,p_z,0),
    (c,s,1)
  \bigr).
\end{equation}
The complete primitive-cell basis is the exact concatenation
\begin{equation}
  \mathcal B
  =
  \mathcal B_{c_0}\mathbin\Vert
  \mathcal B_{c_1}\mathbin\Vert\cdots\mathbin\Vert
  \mathcal B_{c_{N_c-1}},
  \qquad N_b=5N_c.
\end{equation}
For diamond-structure silicon, $N_c=2$ and $N_b=10$.

The symbol $s^{\ast}$ denotes the second radial channel of the $s$ component:
\begin{equation}
  s\equiv(s,0),
  \qquad
  s^{\ast}\equiv(s,1).
\end{equation}
The star is not a principal quantum number, an excited eigenstate label, or a
claim that one continuum atomic orbital has been constructed uniquely.

Basis compatibility is exact and order sensitive. Equal dimensions or equal
orbital counts are insufficient. Two represented matrices may be subtracted
only after the center definitions, component labels, radial channels, ordering,
spin convention, geometry, units, gauge, and energy reference have been
identified.

The present class is orthogonal and spinless. Its overlap matrix is therefore
represented as the identity, and spin--orbit coupling is absent. A
nonorthogonal basis, generalized eigenproblem, explicit spin, or spinor model
is a different model class rather than an optional formatting of this one.

\section{Ordinary-cell hopping representation}
\index{cell gauge}
\index{hopping coefficient}

Let $H_{\mathbf R}\in\mathbb C^{N_b\times N_b}$ denote the coefficient attached
to ordinary-cell displacement $\mathbf R\in\mathbb Z^3$. The cell-gauge Bloch
matrix is defined by the positive-exponent convention
\begin{equation}
  H_{\mathrm{cell}}(\mathbf q)
  =
  \sum_{\mathbf R\in\mathbb Z^3}
  H_{\mathbf R}
  \exp\!\left(+2\pi i\,\mathbf q\cdot\mathbf R\right),
  \label{eq:sk-cell-bloch}
\end{equation}
where $\mathbf q$ is a fractional reciprocal coordinate. Hermiticity of the
Bloch matrix follows from
\begin{equation}
  H_{-\mathbf R}=H_{\mathbf R}^{\dagger}.
  \label{eq:sk-hermitian-support}
\end{equation}

Equation~\eqref{eq:sk-cell-bloch} is a representation convention. Reversing the
Fourier sign changes the definitions of the coefficient matrices and every
subsequent gauge relation. The sign must therefore be retained as part of the
model identity.

The coefficient $H_{\mathbf0}$ may contain both onsite blocks and hopping
between distinct centers inside the same ordinary cell. These are different
scientific roles. In particular, a two-center bond construction does not by
itself define onsite orbital energies.

\section{Orbital-center embedding and reciprocal sewing}
\index{gauge!orbital-center}
\index{reciprocal sewing!tight-binding}

Ordinary-cell displacement and fractional orbital-center embedding are kept
separate. Let $\boldsymbol\tau_a$ be the center coordinate associated with
basis coordinate $a$, and define
\begin{equation}
  D(\mathbf q)
  =
  \operatorname{diag}\!
  \left(
    \exp(+2\pi i\,\mathbf q\cdot\boldsymbol\tau_0),
    \ldots,
    \exp(+2\pi i\,\mathbf q\cdot\boldsymbol\tau_{N_b-1})
  \right).
  \label{eq:sk-center-gauge}
\end{equation}
The full-center representation is
\begin{equation}
  H_{\mathrm{full}}(\mathbf q)
  =
  D(\mathbf q)^{\dagger}
  H_{\mathrm{cell}}(\mathbf q)
  D(\mathbf q).
  \label{eq:sk-full-center}
\end{equation}
Thus a matrix element associated with coefficient displacement $\mathbf R$ and
basis centers $a,b$ carries the total phase
\begin{equation}
  \exp\!\left[
    +2\pi i\,\mathbf q\cdot
    (\mathbf R+\boldsymbol\tau_b-\boldsymbol\tau_a)
  \right].
\end{equation}
Folding $\boldsymbol\tau_b-\boldsymbol\tau_a$ into the integer displacement
would erase the distinction between ordinary-cell translation and intracell
embedding.

The transformation in Eq.~\eqref{eq:sk-full-center} is unitary, so the cell and
full-center matrices have the same spectrum at each $\mathbf q$. Their entries
are not generally equal. For an integer reciprocal shift
$\mathbf G\in\mathbb Z^3$,
\begin{equation}
  D(\mathbf q+\mathbf G)=D(\mathbf q)S_{\mathbf G},
  \qquad
  (S_{\mathbf G})_{aa}
  =\exp(+2\pi i\,\mathbf G\cdot\boldsymbol\tau_a).
\end{equation}
If $H_{\mathrm{cell}}$ is reciprocal periodic, then
\begin{equation}
  H_{\mathrm{full}}(\mathbf q+\mathbf G)
  =
  S_{\mathbf G}^{\dagger}
  H_{\mathrm{full}}(\mathbf q)
  S_{\mathbf G}.
\end{equation}
The full-center matrix is therefore reciprocal covariant rather than
entrywise periodic. Its spectrum remains periodic. This sewing relation must
not be confused with a change in the physical operator.

\section{Directed bonds and physical geometry}
\index{bond!directed}
\index{direction cosine}

A directed bond representative is
\begin{equation}
  b=(\mathbf R,s,t),
  \qquad
  \mathbf R\in\mathbb Z^3,
\end{equation}
with explicit source center $s$ and target center $t$. Its Hermitian reverse is
\begin{equation}
  b^{\mathrm{rev}}=(-\mathbf R,t,s).
\end{equation}
At fixed geometry, the Cartesian displacement from the source center in cell
$\mathbf0$ to the target center in cell $\mathbf R$ is
\begin{equation}
  \Delta\mathbf r_b
  =
  A(\mathbf R+\boldsymbol\tau_t-\boldsymbol\tau_s).
  \label{eq:sk-bond-displacement}
\end{equation}
A two-center hopping bond requires
$\lVert\Delta\mathbf r_b\rVert_2>0$. Its normalized direction
\begin{equation}
  \mathbf d_b=(l_b,m_b,n_b)^{\mathsf T}
  =
  \frac{\Delta\mathbf r_b}
       {\lVert\Delta\mathbf r_b\rVert_2}
  \label{eq:sk-bond-direction}
\end{equation}
provides the direction cosines entering the Slater--Koster angular factors.

The integer displacement $\mathbf R$, the physical displacement
$\Delta\mathbf r_b$, and the normalized direction $\mathbf d_b$ have different
roles. The first controls the cell-gauge Bloch phase, the second carries length
and geometry, and the third is nondimensional. None should be inferred from a
filename, shell label, or matrix index.

A bond inventory must be crystallographically complete for the model class and
must not contain both a representative and the reverse that will already be
generated by Hermiticity. Uniqueness prevents double counting; it does not prove
that all physically relevant interactions have been retained.

\section{Homopolar two-center \texorpdfstring{$sp^3s^{\ast}$}{sp3s*} matrix}
\index{Slater--Koster integral}
\index{two-center approximation}

For one homopolar directed bond with normalized direction
$\mathbf d=(l,m,n)^{\mathsf T}$, define the integral vector
\begin{equation}
  \mathbf V
  =
  \bigl(
  V_{ss\sigma},
  V_{sp\sigma},
  V_{pp\sigma},
  V_{pp\pi},
  V_{s^{\ast}p\sigma},
  V_{ss^{\ast}\sigma},
  V_{s^{\ast}s^{\ast}\sigma}
  \bigr).
\end{equation}
In the ordered basis $(s,p_x,p_y,p_z,s^{\ast})$, the represented two-center
matrix is
\begin{equation}
  T(\mathbf d;\mathbf V)
  =
  \begin{pmatrix}
    V_{ss\sigma}
      & V_{sp\sigma}\mathbf d^{\mathsf T}
      & V_{ss^{\ast}\sigma} \\
    -V_{sp\sigma}\mathbf d
      & V_{pp\pi}I_3
        +(V_{pp\sigma}-V_{pp\pi})
          \mathbf d\mathbf d^{\mathsf T}
      & -V_{s^{\ast}p\sigma}\mathbf d \\
    V_{ss^{\ast}\sigma}
      & V_{s^{\ast}p\sigma}\mathbf d^{\mathsf T}
      & V_{s^{\ast}s^{\ast}\sigma}
  \end{pmatrix}.
  \label{eq:sk-sp3sstar-matrix}
\end{equation}
The signs of the $s$--$p$ and $s^{\ast}$--$p$ blocks are part of the directed
bond convention. They are not removable typographical choices.

Equivalently,
\begin{equation}
  T(\mathbf d;\mathbf V)
  =
  \sum_{\alpha=1}^{7}V_{\alpha}B_{\alpha}(\mathbf d),
  \label{eq:sk-channel-expansion}
\end{equation}
where the seven channel matrices contain only the angular factors. At fixed
geometry and bond inventory, Eq.~\eqref{eq:sk-channel-expansion} is linear in
the two-center integrals. This linearity defines a parameterization; it does
not choose a fitting objective, regularizer, sample weighting, or acceptance
threshold.

For real homopolar integrals,
\begin{equation}
  T(-\mathbf d;\mathbf V)
  =
  T(\mathbf d;\mathbf V)^{\mathsf T}.
  \label{eq:sk-reversal}
\end{equation}
This reversal identity underlies Hermitian cell assembly. It does not extend
without modification to every heteropolar, magnetic, spinorial, or
complex-phase model.

\section{Assembly of ordinary-cell coefficients}
\index{assembly!tight-binding}

Let $E_c\in\{0,1\}^{N_b\times5}$ embed the ordered five-orbital block of center
$c$ into the complete basis, with
\begin{equation}
  E_c^{\mathsf T}E_c=I_5.
\end{equation}
Let each directed representative
\begin{equation}
  b=(\mathbf R_b,s_b,t_b,\mathbf d_b,g_b)
\end{equation}
carry a parameter-group label $g_b$. The label states explicit parameter
sharing only; it does not infer a material, neighbor shell, or chemical meaning.
If $\mathbf V_{g_b}$ is the corresponding integral vector, then assembly adds
\begin{equation}
  H_{\mathbf R_b}
  \mathrel{+}=
  E_{s_b}
  T(\mathbf d_b;\mathbf V_{g_b})
  E_{t_b}^{\mathsf T}
  \label{eq:sk-forward-assembly}
\end{equation}
and generates the reverse contribution
\begin{equation}
  H_{-\mathbf R_b}
  \mathrel{+}=
  E_{t_b}
  T(\mathbf d_b;\mathbf V_{g_b})^{\mathsf T}
  E_{s_b}^{\mathsf T}.
  \label{eq:sk-reverse-assembly}
\end{equation}
When $\mathbf R_b=\mathbf0$ and $s_b\neq t_b$, both directed blocks enter the
same coefficient $H_{\mathbf0}$. A zero-displacement self bond is not a
replacement for an onsite operator.

For a homopolar crystal with cubic onsite symmetry, one common onsite block is
\begin{equation}
  O_c
  =
  \operatorname{diag}
  (\varepsilon_s,\varepsilon_p,\varepsilon_p,\varepsilon_p,
   \varepsilon_{s^{\ast}}),
  \label{eq:sk-onsite}
\end{equation}
embedded as $E_cO_cE_c^{\mathsf T}$ in $H_{\mathbf0}$. The exact onsite class,
energy reference, and sharing across centers remain independently declared
model choices.

Equations~\eqref{eq:sk-forward-assembly} and
\eqref{eq:sk-reverse-assembly} produce an ordinary-cell representation. The
fractional center phases are applied only afterward through
Eq.~\eqref{eq:sk-full-center}. Neighbor generation, distance scaling, onsite
selection, fitting, and material admission are not hidden inside assembly.

\section{Diamond-silicon model classes}
\index{silicon!Slater--Koster model}
\index{model hierarchy!tight-binding}

For diamond silicon, the primitive cell contains two equivalent centers. The
spinless orthogonal spatial basis therefore has ten coordinates. In conventional
cubic Cartesian axes with lattice constant $a$, the four nearest-neighbor unit
directions from one sublattice may be written
\begin{equation}
  \frac{1}{\sqrt3}
  \bigl
  \{
    (1,1,1),
    (1,-1,-1),
    (-1,1,-1),
    (-1,-1,1)
  \bigr\},
  \label{eq:diamond-nearest-directions}
\end{equation}
with bond length $a\sqrt3/4$. They connect opposite sublattices. The twelve
second-neighbor directions lie in the $\langle110\rangle$ family, have length
$a/\sqrt2$, and connect sites on the same fcc sublattice. These Cartesian shell
descriptions do not replace the exact ordinary-cell displacement and center
labels required by Eq.~\eqref{eq:sk-bond-displacement}.

A concrete candidate class must additionally freeze

\begin{itemize}
  \item the primitive lattice and two center positions;
  \item the exact ten-coordinate order;
  \item the onsite parameter inventory;
  \item the first-, second-, and any later-neighbor bond inventories;
  \item which bonds share each two-center integral vector;
  \item the cell/full-center gauge convention and reciprocal sewing;
  \item the energy unit and scalar energy reference; and
  \item the parameter domain and any sign or symmetry restrictions.
\end{itemize}

A nearest-neighbor class and a second-neighbor class are different model
classes even when both use the same ten-coordinate basis. Adding a shell changes
support and parameter sharing; it is not merely refining a numerical solver.
Likewise, changing the basis, allowing overlap, adding spin--orbit coupling, or
introducing a momentum-dependent alignment changes the scientific model class.

The two-center construction does not assert that a nearest-neighbor
$sp^3s^{\ast}$ class is adequate for silicon. Adequacy requires the spectral,
aligned-operator, withheld-observable, and error-accounting tests defined in
later chapters.

\section{Common-space admission before operator comparison}
\index{common space!tight-binding and Wannier}
\index{principal angles!trial and retained spaces}

A Wannier Hamiltonian and a Slater--Koster Hamiltonian are initially represented
on different coordinate spaces. Equal rank does not identify those spaces. Let
$W(\mathbf q)$ be a retained parent-to-Wannier frame, let $A(\mathbf q)$ contain
trial-orbital projections in the parent space, and let $T$ map the declared
trial order to the canonical tight-binding order. A projection-anchored overlap
has the form
\begin{equation}
  B(\mathbf q)
  =
  W(\mathbf q)^{\dagger}A(\mathbf q)T^{\dagger}.
  \label{eq:sk-common-space-overlap}
\end{equation}
A rank-$N_b$ unitary identification obtained by L\"owdin
orthonormalization requires
\begin{equation}
  \operatorname{rank}B(\mathbf q)=N_b
  \quad\text{for every admitted }\mathbf q,
  \label{eq:sk-full-rank-gate}
\end{equation}
or, numerically, a smallest singular value resolved above the declared source
and arithmetic uncertainty. If this gate fails, no unitary common-space map is
defined under the stated projection contract, and an aligned operator residual
must not be formed by singular-value flooring or an undeclared pseudoinverse.

The singular values of $B$ are invariant under unitary changes of frame inside
the retained subspace. If
\begin{equation}
  W(\mathbf q)\longmapsto W(\mathbf q)V(\mathbf q),
  \qquad V(\mathbf q)\in U(N_b),
\end{equation}
then
\begin{equation}
  B(\mathbf q)\longmapsto V(\mathbf q)^{\dagger}B(\mathbf q),
\end{equation}
which leaves all principal angles unchanged. Energy degeneracy alone therefore
does not repair a deficient trial-to-retained overlap. A rank failure is a
state-space admission result; it is not yet a fit result or proof that the
Slater--Koster model class is spectrally inadequate.

A different successor may instead posit one momentum-independent global unitary
$Q\in U(N_b)$ and compare
\begin{equation}
  H_{\mathrm W}(\mathbf q)
  \quad\text{with}\quad
  QH_{\mathrm{TB}}(\mathbf q;\boldsymbol\theta)Q^{\dagger}.
\end{equation}
That defines a different alignment family. Allowing an arbitrary
$Q(\mathbf q)$ would generally absorb momentum-dependent gauge structure and
could alter real-space locality, so it is not equivalent to one global orbital
identification.

\section{Scientific scope and exclusions}
\index{model boundary!Slater--Koster}

The model defined here supports statements about a finite orthogonal,
spinless, homopolar two-center representation with explicit centers and a
specified bond inventory. It does not by itself support claims about

\begin{itemize}
  \item heteropolar bond rules or element-dependent directed integrals;
  \item overlap matrices or generalized eigenproblems;
  \item explicit spin, magnetism, or spin--orbit coupling;
  \item $d$ or $f$ orbitals;
  \item environment-dependent, many-center, or self-consistent interactions;
  \item strain scaling without a separately defined geometric law;
  \item optical, velocity, or other response matrix elements;
  \item parameter transferability between geometries or materials; or
  \item scientific adequacy for silicon.
\end{itemize}

Those extensions may be scientifically important, but each changes the model
class or its validation requirements. The role of the present chapter is to
make the admitted operator, basis, gauge, and two-center assumptions explicit
before reduction and model selection begin.
